\documentclass{article}
\usepackage[nodayofweek]{datetime}
\usepackage{graphicx}
\usepackage{float}

\title{Research Proposal: \\
        Optimising PSMA PET Segmentation using the Mamba Architecture}
\author{Author: James Wigfield (23334375)\\ 
        Supervisor: Dr. Jake Kendrick \\ 
        Co-supervisor: Dr. Mubashar Hassan \\ 
        Degree: Bachelor of Advanced Computer Science (Honours)}
\date{\monthname[\the\month] \the\year}

% for creation of the timeline PDF from the mermaid markdown file
%\immediate\write18{mmdc -i timeline.mmd -o timeline.pdf}

\begin{document}

\maketitle

\section{Introduction and Background}
\subsection{Clinical Context}
Prostate cancer (PCa) is one of the most common cancers amongst men globally \cite{Bray2024GlobalCS}, however, 
localised PCa is generally associated with high survival rates. Biochemical recurrence (BCR) following treatment
frequently leads to metastatic prostate cancer (mPCa) which is why accurately detecting and localising mPCa 
lesions is critical for guiding treatment decisions and improving patient outcomes. $[^{68}Ga]$Ga-PSMA-11 Positron 
Emission Tomography (PET) allows clinicians to visualise and locate mPCa lesions throughout the body. However, 
analysing and segmenting these whole-body scans manually carries a significant workload and is subject to inter-observer
variability. Therefore, the need for a fully automated segmentation tool is evident.

\subsection{Current Baseline}
The current baseline for PSMA PET segmentation is the nnU-Net architecture which has achieved some very promising
results in the field of medical image segmentation. This baseline model achieves a 94.5\% patient-level accuracy (121/128),
and for lesion-level detection achieved an F1 score of 79.9\% which balances a high positive predictive value (PPV) of 88.2\% 
against a slightly lower sensitivity of 73\% \cite{kendrick2022fully}. Additionally, applying a cascading 3D CNN to whole-body
PSMA PET scans demands significant computational resources and latency during inference. The proposed 3D Mamba architecture 
will aim to improve this F1 score while accelerating clinical inference times.


\subsection{Computational Limits}
Deep 3D convolutional neural networks (CNNs) face severe computational bottlenecks due to GPU memory constraints, forcing 
practitioners to make trade-offs between patch size and batch size. This limitation becomes particularly critical when processing 
whole-body imaging scans, where volumetric data exceeds typical GPU memory allocations. 

\subsection{State Space Models (SSMs) and Mamba}
3D CNNs such as the nnU-Net baseline achieve robust volumetric segmentation leveraging hierarchical encoder-decoder architectures. 
However, 3D convolutions are computationally expensive, which forces restricted patch sizes that limit the model's global spatial 
context \cite{isensee2021nnu}. Using Transformer-based models to capture these missing long-range dependencies is not viable due to 
their quadratic scaling with respect to sequence length. To overcome these limitations, State Space Models (SSMs) present a promising
alternative which scale linearly with sequence length. The Mamba architecture introduces a selection mechanism where the SSM parameters
are generated based on the input, allowing the model to selectively filter out irrelevant information (such as healthy background tissue) 
and remember relevant small features (such as small lesions) \cite{gu2023mamba}.

\section{Problem Statement and Hypothesis}
By leveraging the linear complexity and selective memory of State Space Models (SSMs), a 3D Mamba
architecture can process whole-body PSMA PET scans with greater computational efficiency than the 
existing baseline models. Additionally, optimising the model for lesion-level detection will result in
a clinically actionable tool that can accurately identify metastatic lesions. 

\section{Project Aims}
The aims of this research project are as follows:

\begin{enumerate}
    \item Design a 3D Mamba-based architecture to specifically address the extreme volumetric disparity
    between tiny metastatic lesions and vast healthy background tissue in PSMA PET scans.
    \item Benchmark the proposed architecture against the rigorously configured nnU-Net baseline using 
    a PSMA PET dataset. 
    \item Evaluate using a dual-metric approach, prioritising lesion-level F1 score to ensure an optimal
    balance between sensitivity and precision (PPV), alongside voxel-level Dice Similarity Coefficient (DSC) 
    to ensure critical tumour detection.   
\end{enumerate}

% again, do not focus on sensitivity alone. also precision. F1 score should capture this.

\section{Methodology}

\subsection{Dataset and Preprocessing}
The model will be trained and evaluated on the Sir Charles Gairdner Hospital (SCGH) dataset, which consists of whole-body PSMA PET 
scans of around 300 patients \cite{mccarthy2019multicenter}. In addition to this dataset, an open source autoPET dataset will be used along with possibly other open
source datasets to further augment the training data. 


% there's another dataset that joel has, but mentioning SCGH dataset and opensource should be sufficient.

\subsection{Model Architecture Design}
The proposed architecture will use a hybrid CNN-Mamba design drawing inspiration from frameworks like U-Mamba \cite{ma2024umamba}. 
This design will leverage the capabilities of convolutional layers for local feature extraction with the long-range dependency
modelling capabilities of State Space Models (SSMs) to capture global context across the entire scan.

The encoder will consist of convolutional residual blocks to capture local features. The resulting spatial features will be flattened
into 1D sequences so the 3D volumetric PET scans can be processed by the Mamba blocks \cite{ma2024umamba}.

To mitigate the loss of high-resolution spatial information, the architecture will incorporate skip connections between the encoder
and decoder. 

While the initial network will be based on a hybrid CNN-SSM design, experiments will also be conducted with pure 3D mamba blocks versus a hybrid
CNN-Mamba design to evaluate the trade-offs between computational efficiency and segmentation performance.

\subsection{Volumetric VRAM Mitigation}
To tackle the GPU memory constraints that limit 3D CNN patch sizes, this methodology will explore Mamba's hardware-aware parallel scan
algorithm. This algorithm allows the model to compute the sequence recurrence directly in SRAM and avoids materialising the entire
sequence in VRAM, which drastically reduces the memory footprint and allows for processing larger volumetric data without sacrificing
batch size \cite{gu2023mamba}.

Since Mamba scales linearly, the training pipeline will be designed to process significantly larger volumetric patches than the baseline 
nnU-Net. This will ensure the model retains vast spatial context across the whole-body scans which is required to map relationships between
a primary prostate tumour and its metastatic lesions.

In order to generate full-body segmentations during inference without triggering out-of-memory (OOM) errors, an overlap-tile approach will
be implemented \cite{ronneberger2015unet} but with a 3D volumetric adaptation. To ensure prediction accuracy at the patch borders, the 
model will extrapolate missing boundary context via mirror padding. 

\subsection{Loss Function Formulation}
In whole-body PSMA PET scans, healthy background voxels vastly outnumber the voxels of small metastatic lesions. Training with standard
loss functions in this situation can lead to models that are biased towards high precision but low sensitivity. To mitigate this 
imbalance, the training pipeline will investigate alternative loss functions such as the Tversky loss which is a generalisation of the 
DSC and $F_\beta$ scores. The Tversky index introduces hyperparameters $\alpha$ and $\beta$ to control the trade-off between false positives 
and false negatives \cite{salehi2017tversky}.

In oncology, false negatives (missing a metastatic lesion) are generally more detrimental than false positives. Setting a higher $\beta$ value 
in the Tversky loss will penalise false negatives more heavily, encouraging the model to prioritise sensitivity while still maintaining a 
reasonable precision. This approach will directly target the 73\% sensitivity of the baseline nnU-Net model while maintaining or improving 
the high PPV of 88.2\% \cite{kendrick2022fully}.

While the Tversky loss will be the primary focus, the research will also experiment with alternative loss functions such as focal loss and dice 
loss to evaluate their impact on the model's performance. 

\section{Resources and Hardware}
Model training will require a GPU with CUDA support with at least 16GB of VRAM. Primary training will be conducted on the SCGH Medical 
Physics Department's High Performance Computing machines. 

\vspace{0cm}
\section{Timeline}
\begin{figure}[H]
    \centering
    % It grabs the newly generated PDF and drops it in
    \includegraphics[width=\textwidth,scale=1.2]{timeline.pdf}
    %\caption{Project Timeline Gantt Chart}
\end{figure}
\vspace{0cm}

\section{Statement of AI Use}
The following AI tools were used in the preparation of this research proposal:

\begin{itemize}
    \item \textbf{Google Gemini 3.1} - Provided a rough structural outline of the research proposal and answered terminology questions relating to referenced research papers.
    \item \textbf{Claude Haiku 4.5 (GitHub Copilot)} - Assisted with spelling checks, converting technical terms and mathematical parameters into LaTeX symbols, and general LaTeX formatting guidance.
\end{itemize}
 

\bibliography{references}
\bibliographystyle{plain}
\end{document}

